\documentclass[border=8pt]{standalone}
\usepackage{fontspec}
\setmainfont{Noto Sans}[Ligatures=TeXOff]
\XeTeXgenerateactualtext=1
\setsansfont{Noto Sans}
\newfontfamily\bengalifont{Noto Sans Bengali}[Script=Bengali, NFSSFamily=bnfam]
\usepackage[Latin,Bengali]{ucharclasses}
% group-free font switching (as in header-common.tex): robust inside braces, cells and quotes
\makeatletter
\def\BN@saved{\rmdefault}
\def\BN@bn{bnfam}
\newcommand\BNin{\edef\BN@tmp{\f@family}\ifx\BN@tmp\BN@bn\else\let\BN@saved\BN@tmp\fi\fontfamily{bnfam}\selectfont}
\newcommand\BNout{\fontfamily{\BN@saved}\selectfont}
\makeatother
\setTransitionsFor{Bengali}{\BNin}{\BNout}
% Typographic quotes and dashes are not in a ucharclasses class (class 0), so a Bangla word right
% after “ or ‘ would not trigger the Bangla font. Treat them as Latin so that every switch fires.
\XeTeXcharclass"2018=\XeTeXcharclass`A
\XeTeXcharclass"2019=\XeTeXcharclass`A
\XeTeXcharclass"201C=\XeTeXcharclass`A
\XeTeXcharclass"201D=\XeTeXcharclass`A
\XeTeXcharclass"2013=\XeTeXcharclass`A
\XeTeXcharclass"2014=\XeTeXcharclass`A
\XeTeXcharclass"2026=\XeTeXcharclass`A
\usepackage{newunicodechar}
\newunicodechar{।}{{\bengalifont\char"0964}}
\usepackage{tikz}
\usetikzlibrary{arrows.meta,positioning,shapes.geometric,shapes.misc,shapes.symbols,fit,backgrounds,calc,chains}
\definecolor{ink}{HTML}{111827}
\definecolor{band}{HTML}{F3F4F6}
\definecolor{band2}{HTML}{E5E7EB}
\definecolor{accent}{HTML}{1D4ED8}
\definecolor{accentlight}{HTML}{DBEAFE}
\definecolor{warn}{HTML}{B91C1C}
\definecolor{warnlight}{HTML}{FEE2E2}
\definecolor{ok}{HTML}{047857}
\definecolor{oklight}{HTML}{D1FAE5}
\tikzset{
  >={Stealth[length=2.6mm]},
  box/.style={draw=ink, line width=0.6pt, rounded corners=9pt, fill=white, align=center,
              inner sep=4pt, font=\small},
  pill/.style={box, rounded corners=13pt},
  key/.style={box, fill=accentlight, draw=accent, line width=0.9pt},
  soft/.style={box, fill=band},
  dec/.style={draw=ink, line width=0.6pt, diamond, aspect=2.2, align=center, inner sep=1.5pt,
              font=\footnotesize, fill=white},
  arr/.style={->, draw=ink, line width=0.6pt},
  darr/.style={->, draw=ink, line width=0.6pt, dashed},
  lbl/.style={font=\footnotesize, fill=white, inner sep=1.5pt},
  layer/.style={draw=band2, fill=band, rounded corners=6pt},
  title/.style={font=\bfseries\small, align=left, anchor=west},
}
\begin{document}
\sffamily\hyphenpenalty=10000\exhyphenpenalty=10000

\begin{tikzpicture}
\tikzset{stage/.style={box, minimum width=96mm, minimum height=17mm, text width=92mm, font=\footnotesize},
         gate/.style={dec, aspect=3.2, text width=44mm, font=\footnotesize}}
\node[key, stage] (s0) at (0,0) {\textbf{পর্যায় ০} \quad মাস ১ থেকে ১২\\[1mm]বিদ্যমান উন্মুক্ত মডেলে উদ্ধারভিত্তিক ব্যবস্থা; কোনো প্রশিক্ষণ নয়; BanglaPolicyQA v0};
\node[gate] (g1) at (0,-25mm) {ঘাটতি সীমার ওপরে, ডেটা ও কম্পিউটিং প্রস্তুত?};
\node[stage] (s1) at (0,-50mm) {\textbf{পর্যায় ১} \quad মাস ১৩ থেকে ৩০\\[1mm]নিয়ন্ত্রিত উন্মুক্ত কর্পাসে ডোমেইন অভিযোজন (LoRA); নির্দেশনা-টিউনিং; মডেল কার্ড};
\node[gate] (g2) at (0,-75mm) {উল্লেখযোগ্য উন্নতি ও নিরাপত্তা পরীক্ষা উত্তীর্ণ?};
\node[stage] (s2) at (0,-100mm) {\textbf{পর্যায় ২} \quad মাস ৩০-এর পর, শুধু যৌক্তিক হলে\\[1mm]যৌথ জাতীয় কম্পিউটিংয়ে অব্যাহত প্রাক-প্রশিক্ষণ; স্বাধীন নিরীক্ষা};
\draw[arr] (s0) -- (g1);
\draw[arr] (g1) -- node[lbl, right]{হ্যাঁ} (s1);
\draw[arr] (s1) -- (g2);
\draw[arr] (g2) -- node[lbl, right]{হ্যাঁ} (s2);
\node[box, dashed, font=\footnotesize, text width=34mm] (st1) at (74mm,-25mm) {না: বর্তমান পর্যায়ে থাকা};
\node[box, dashed, font=\footnotesize, text width=34mm] (st2) at (74mm,-75mm) {না: বর্তমান পর্যায়ে থাকা};
\draw[darr] (g1) -- (st1); \draw[darr] (g2) -- (st2);
\node[soft, minimum width=150mm, text width=146mm, minimum height=9mm, font=\footnotesize] at (14mm,-120mm) {প্রতিটি পর্যায়: উন্মুক্ত বেঞ্চমার্ক, মূল্যায়ন প্রতিবেদন, ERRC পর্যালোচনা, DOI-সহ সংরক্ষিত রিলিজ};
\end{tikzpicture}
\end{document}
